% !TeX root = main.tex

\section{Software Improvements}
\label{sec:Software-Improvements}

This section describes the software engineering efforts of the last two years. Our goal is to improve and sustain the scientific community efforts around QMCPACK's evolution in the current and future landscape of HPC.

% Michael
\subsection{Docker containers}

Containers are a lightweight alternative to virtual machines to create portable and reproducible environments without a full operating system.
Due to its nature as a complex physical simulation framework, QMCPACK has several configuration options to build the final executable.
We use pre-packaged Docker images to provide uniformity across environments used in GitHub Actions CI tests, including appropriate dependencies hosted in DockerHub.
In addition, ``Dockerfiles'' used for image configuration are available in the QMCPACK repository for reproducibility.
These containers can be pulled to any Linux host to trigger an interactive debug session in the same environment used for CI.
On the other hand, we do not use containers for macOS and GPU runs as there is currently limited value due to the small number of targeted use cases and the bare-metal nature of testing on GPUs. 

In addition, we have begun exploring using the Spack package manager to generate Docker images and leverage its package managing capabilities.
This integration and interoperability simplify the distinct processes each dependency may require to build and install correctly on different operating systems and environments.

% Philip: narrative, look at ECP report for baseline
\subsection{Strategic CI expansion on GitHub Actions}

Starting in mid-2021, QMCPACK began the adoption of GitHub Actions in place of Jenkins for CI.
The main reasons are:
(i) making full utilization of the freely provided GitHub CPU runners in virtual machines (VMs) using our Docker containers, and
(ii) providing CI for GPU development via runners hosted at ORNL equipped with NVIDIA and AMD cards. In addition, GitHub Actions offers simple workflow configurations via YAML files, provides reusable automated functionality via open-source workflow ``actions'' and other services (e.g., Codecov, DockerHub, GitHub Actions data encryption via ``secrets''), and enables data mining access of the logs via the GitHub command line interface (CLI). This adoption allowed for the desired sustainability goals of lowering maintenance costs while adding relevant validation and verification checks to QMCPACK.

QMCPACK uses GitHub Actions VMs to perform an automatic initial check for CI using Docker containers on available virtual machines (VMs) for the configurations listed in Table~\ref{tab:GitHubActions}.
While the main target is Linux, macOS testing using the Accelerate framework for linear algebra is also provided.
The current set of CI configurations has been evolving to prioritize checks that address problematic and error-prone areas of the code.
As such, address (ASan) and undefined behavior (UBSan) sanitizer checks and code coverage using GNU's gcov were added early in the CI migration, while jobs running in serial mode (NoOMP, NoMPI) were added later. 
The latter ensures the code can be built with the fewest number of dependencies and in the easiest mode to utilize debuggers and profilers.
An important benefit is that offloading the hardware maintenance and system administration aspects lowers the overall cost of adding more configurations.
The trade-off is that while each job runs on limited hardware resources at small scales, they are all triggered concurrently when new contributions are opened via pull requests targeting the main development branch. The configurations have been curated to keep turnaround in the CI at close to one hour of real-time.
 
\begin{table}[!ht]
\centering
\caption{CI on GitHub Actions hosted systems (CPU-only) as of March 2023}
\begin{tabular}{||p{0.35\linewidth} | p{0.55\linewidth}||}
\hline
System VM     & Job Configuration  \\
\hline \hline 
\textit{linux}          & GCC9-NoMPI-Debug-Real         \\
2-core CPU (x86\_64)    & GCC9-NoMPI-NoOMP-Real           \\
7GB RAM, 14GB SSD       & GCC9-NoMPI-NoOMP-Complex        \\
                        & GCC9-NoMPI-Sandbox-Real         \\
                        & GCC9-MPI-Gcov-Real              \\
                        & GCC9-MPI-Gcov-Complex           \\
                        & GCC11-NoMPI-Werror-Real       \\
                        & GCC11-NoMPI-Werror-Complex       \\
                        & GCC11-NoMPI-Werror-Real-Mixed       \\
                        & GCC11-NoMPI-Werror-Complex-Mixed       \\
                        & Clang10-NoMPI-ASan-Real       \\
                        & Clang10-NoMPI-ASan-Complex       \\
                        & Clang10-NoMPI-UBSan-Real      \\
                        & Clang12-NoMPI-Offload-Real       \\
\hline
\textit{macOS}          & GCC11-NoMPI-Real    \\
3-core CPU (x86\_64)    & \\
14GB RAM, 14GB SSD      & \\
\hline
\end{tabular}
\label{tab:GitHubActions}
\end{table}

Testing each pull request requires running the proposed changes on GPUs.
This is achieved via a second-level CI workflow, which must be triggered manually (by authorized members).
The specifications for each system are provided in Table~\ref{tab:self-hosted} along with the job configurations run on each system (sulfur and nitrogen). At the time of this work the current systems had a Red Hat Enterprise Linux (RHEL) operating system version 8.
This strategy meets our security requirements to prevent access to unauthorized agents, thus adding another layer of security on top of GitHub's two-factor authentication and branch
protections.
It is worth noting that the CPU and GPU testing runs concurrently.
Due to limited GPU memory, GPU jobs are, however, run sequentially.
To provide a partial test of the OpenMP target offload GPU implementation  in the first tier of CI we utilize the ability of OpenMP offload code to be compiled by LLVM for host CPUs.
This provides an initial screen for modifications to the GPU code without needing to commit any GPU hardware.

\begin{table}[!h]
\centering
\caption{CI on ORNL self-hosted systems as of March 2023}
\begin{tabular}{||p{0.29\linewidth} | p{0.58\linewidth}||}
\hline
System Specs      & Job Configuration  \\
\hline \hline 
\textit{sulfur Linux RHEL8}     &                                           \\
                    &                                           \\
\textbf{CPU}        & GCC8-NoMPI-MKL-Real-Mixed                 \\
2xIntel Xeon        & GCC8-NoMPI-MKL-Complex-Mixed              \\
Gold 6248R 24-core  & GCC8-NoMPI-MKL-Real                       \\
                    & GCC8-NoMPI-MKL-Complex                    \\
                    &                                           \\
\textbf{GPU}        &                                           \\
NVIDIA Tesla V100   & Clang15-MPI-CUDA-AFQMC-Offload-Real-Mixed \\
                    & Clang15-MPI-CUDA-AFQMC-Offload-Real       \\
                    & Clang15-MPI-CUDA-AFQMC-Offload-Complex-Mixed  \\
                    & Clang15-MPI-CUDA-AFQMC-Offload-Complex    \\
                    & Intel19-MPI-CUDA-AFQMC-Real-Mixed         \\
                    & Intel19-MPI-CUDA-AFQMC-Complex-Mixed      \\
                    & Intel19-MPI-CUDA-AFQMC-Real               \\
\hline
\textit{nitrogen Linux RHEL8}  &                  \\
2xAMD EPYC          & \\
7542 32-Core        & \\
                    &                                \\
\textbf{GPU}        &                                \\
NVIDIA Tesla V100   & GCC8-MPI-CUDA-AFQMC-Real-Mixed \\
                    & GCC8-MPI-CUDA-AFQMC-Real  \\
                    & GCC8-MPI-CUDA-AFQMC-Complex-Mixed \\
                    & GCC8-MPI-CUDA-AFQMC-Complex  \\
                    &                                           \\
\textbf{GPU}        & \\ 
AMD Vega 20         & ROCm-Clang13-NoMPI-CUDA2HIP-Real-Mixed \\
                    & ROCm-Clang13-NoMPI-CUDA2HIP-Real \\
                    & ROCm-Clang13-NoMPI-CUDA2HIP-Complex-Mixed \\
                    & ROCm-Clang13-NoMPI-CUDA2HIP-Complex \\
\hline
\end{tabular}
\label{tab:self-hosted}
\end{table}

To quantify the impact of the preventive nature of the CI system, we measure the total number of occurrences and the failure rates per month since the beginning of this effort.

\begin{figure}[ht]
    \centering
    \subfloat[GitHub Actions CI.]{
        \includegraphics[width=3.55in]{Figures/GitHubActionsCI_monthly_totals_stacked.pdf}
        \label{fig:totals_gha}
    }
    \vspace{0.1in}
    \subfloat[Self-hosted CI.]{
        \includegraphics[width=3.55in]{Figures/ornlCI_monthly_totals_stacked.pdf}
        \label{fig:totals_ornl}
    }
    \caption{Total workflows.}
    \label{fig:totals}
\end{figure}

Figure~\ref{fig:totals} shows the total number of workflows (a collection of jobs) run in each month, also as an indication of the volume of development activity.
Note the difference in scale (100 \textit{vs.} 300) between the number of workflows run on self-hosted ORNL runners (Fig. \ref{fig:totals_ornl} ``ornl CI'') and those run on the free GitHub-hosted runners (Fig.
\ref{fig:totals_gha} ``GitHub Actions CI'').
This is explained due to our loose requirement for passing GitHub Actions CI first-level checks prior to using the self-hosted runners, which are a more constrained resource.

Failure rates can be extremely variable, and in the case of cloud-based CI in Figure~\ref{fig:totals_gha}, it may also be dependent on external factors (\textit{e.g.}, Codecov failed requests or GitHub Actions incidents). Figure~\ref{fig:totals_ornl} illustrates the importance of having GPU CI, especially around times of high development activity. 

Testing implicitly touches the entire software stack (the Nexus workflow system, various wavefunction conversion utilities, etc), not only in the main QMCPACK application.
The high failure rate experienced in the first quarter of 2023 reflects problems and updates in the test infrastructure.
Peaking in March of 2023, the failures are explained due to the process of upgrading both CI systems without pausing development while issues were resolved.
We introduced Spack-based CentOS CI on GitHub Actions runners, one configuration for which led to the discovery of a persistent bug that was determined to be caused by the OpenBLAS threading policy until a workaround was implemented in QMCPACK and promptly fixed by the OpenBLAS developers.
At the same time, our self-hosted runners had an operating system upgrade (from RHEL8 to RHEL9), which led to several refactoring efforts in sorting out compatibility of CPU and GPU dependency versions. The latter is a good case for how scientific software priorities might differ from other types of development.

\begin{figure}[ht]
    \centering
    \subfloat[``GitHub Actions CI'']{
        \includegraphics[width=3.55in]{Figures/GitHubActionsCI_monthly_failed_jobs.pdf}
        \label{fig:failed_jobs_gha}
    }
    \vspace{0.1in}
    \subfloat[``Self-hosted CI'']{
        \includegraphics[width=3.55in]{Figures/ornlCI_monthly_failed_jobs.pdf}
        \label{fig:failed_jobs_ornl}
    }
    \caption{Failed jobs.}
    \label{fig:failed_jobs}
\end{figure}

To provide information on the nature of the failed CI jobs, Figure~\ref{fig:failed_jobs} shows the total number of failed jobs (not workflows) per month.
In the case for GitHub Actions CI, Figure~\ref{fig:failed_jobs_gha}, the majority of failed jobs come from configurations targeting the Linux runners.
Still, failures caught on the macOS free runners on GitHub Actions are not negligible and provide coverage where many new users initially experiment with the code.
It can be observed in Fig.~\ref{fig:failed_jobs_ornl} that since the beginning of 2022, the jobs identifiers were refactored to meet current standard names, similar to those on the longer-running nightly tests, so they can be classified by job name pointing at a particular configuration type as those shown in Table~\ref{tab:self-hosted}.
Failure counts illustrate the level of development activity on new features, high activity around NVIDIA GPU implementations (sulfur and nitrogen -cuda) early in 2022, while the second half indicates high levels of activity in AMD GPU implementation features (nitrogen-rocm).
GPU testing was able to catch several problems before they made it into the mainline development branch. 
Thus, these checks improve the reliability and robustness of QMCPACK while lowering overall development costs as actively developed GPU technologies targeting exascale systems continue to mature.

% Steve, William
\subsection{Sanitizing critical code paths}

Memory safety has been a long-standing part of technical debt in legacy C and C\texttt{++} code~\cite{10.1145/3129743.3129745}.
Memory leaks, the lack of release of out-of-scope acquired resources, add pressure in the target operating system and threaten long-running simulations.
Through the compile-time instrumentation of address sanitizers~\cite{serebryany2012addresssanitizer}, several memory leaks were identified in QMCPACK unit and functional tests due to the use of raw pointers in pre-C\texttt{++}11 development.
To our knowledge, no memory safety checks were applied to QMCPACK in its early development stages.
We replaced raw pointer instances with ``smart pointers'' introduced in the C\texttt{++}11 standard library, \texttt{std::unique\_ptr} and \texttt{std::shared\_ptr}, to manage object lifetimes automatically following the \emph{resource allocation is initialization} - RAII - paradigm.
The overall process was incremental in nature as each leak provided an opportunity to refactor the affected
object lifetime to favor the use of ``unique'' pointer instances rather than the more permissive ``shared'' pointers.
As each memory leak was resolved, we reinstated each failing test in the GitHub Actions CI (ASan) jobs listed in Table~\ref{tab:GitHubActions} to future-proof critical code execution paths.
We favor the use of sanitizers due to smaller overheads compared to purely run-time approaches such as Valgrind memcheck. 

\begin{table}[h]
\centering
\begin{tabular}{|l|c|}
\hline
Metric type & Occurrences  \\
\hline
leaky unit tests                  & 10   \\
leaky functional tests            & 10   \\
direct leaks                      & 25   \\
indirect leaks                    &  1    \\
\hline 
\end{tabular}
\caption{Metrics showing the scope and nature of the resolved memory leaks in QMCPACK}
\label{tab:leaks}
\end{table}

Table~\ref{tab:leaks} shows the number of failures in the unit and deterministic tests run on CI.
There are two types of leaks to be addressed:
(i) direct, in which allocation and deallocation must be done in QMCPACK, and 
(ii) indirect, in which allocation and deallocation request is done outside QMCPACK, but triggered through the use of public functions in a dependency, e.g. libxml2.
Further details can be found by tracking the issue link in Appendix~\ref{ap1:Artifact}.
The vast majority of these leaks are of a ``direct'' nature.
There were no cases where we had to fall back to shared ownership.
Due to the existence of interoperable C\texttt{++} and C code in QMCPACK, we supplied user-defined custom destructors to \texttt{std::unique\_ptr} to properly free C-style structs.
The only indirect leaks that were resolved were related to the use of C functions in the \texttt{libxml2} library dependency in which the lack of destructor calls was addressed.

% William
\subsection{Code refactoring}

We list some of the important aspects that are currently targeted in QMCPACK's strategic code refactoring efforts that will result in a greater return for the community.

\paragraph{Checkpoint-restart} software running at scale should be able to recover from failure or continue running beyond a maximum allowable time.
QMCPACK uses a checkpoint-restart approach where the state of the simulation is regularly stored to disk. Checkpoint files contain enough data written as plain text, XML and HDF5 files to restart mid-run and are deterministic so as not to affect the results.
Since QMCPACK uses Monte Carlo, this includes restarting of the random number generators.
Part of the effort also includes refactoring I/O components in terms of object ownership and lifetime and improved reusability in the interactions with the underlying HDF5 library.

\paragraph{Phasing out legacy GPU functionality.}
As part of the ECP application development effort, QMCPACK developed a new design and strategy to exploit multiple architectures in a performance-portable manner as described by Luo et al.~\cite{10024599}.
During the development stage of the new performant drivers, users could rely on the original CUDA-specific GPU implementation.
This ``legacy'' implementation was incompatible with CPU-only drivers and resulted in entirely separate calls in the wavefunction calculation phase.
Once the new drivers reached an acceptable level of performance and maturity, the legacy drivers were completely removed, with a reduction of nearly $40\,K$ lines of code.
This illustrates the importance of a software evolution strategy that is able to meet the demands of the science that comes out of new HPC available hardware.
In QMCPACK's case, this strategy is expected to pay off as more exascale and heterogeneous systems continue to be deployed.
By maximizing the number of lines of common code between architectures, maintenance costs are reduced and higher code quality can be expected.

\paragraph{Code coverage via unit tests}
QMCPACK's test suite on GitHub Actions CI uses gcov builds and reports the results to
Codecov to track the progress for code coverage in our tests. 
Pull request contributors will be notified how the new code affects current coverage levels. Decreasing coverage levels will issue a failed CI status.
Code coverage has been increasing since the addition of these CI checks from 38\% (absolute) to 52\% almost linearly with time as shown in Figure~\ref{fig:coverage}.
This coverage (complete line coverage) is expected to keep increasing and already includes all the core functionality of the application. It has been sufficient to catch problems in the greater software stack, although it does not include GPU offload cases. This metric is sensitive to the gcov version, as shown in Figure~\ref{fig:coverage}, the increase in the last two months is due to a GCC toolchain upgrade from version 9 to 11.
While not a specific goal but rather a good practice, progressive migration of code coverage from functional to unit tests is part of our long term goals.
We aim to target the scope of functional tests towards correctness and science cases;
unit tests cover the intended basic software functionalities.

\begin{figure}[ht]
    \centering
    \includegraphics[width=3.55in]{Figures/coverage.pdf}
    \caption{Monthly evolution of QMCPACK code line coverage reported by CI gcov tests.}
    \label{fig:coverage}
\end{figure}

\paragraph{Input validation} QMCPACK's native input is structured in XML files.  While this choice allows for a close relationship between the input data and the underlying modular class framework, the distributed parsing of input on a class by class basis complicates standardization of parsing and validatation of input data.  In concert with the development of performance portable drivers, an improved means of standardized input parsing has been introduced.  In the improved implementation, a fully typed input specification is encoded in association with each underlying class and parsing is deferred to a common base class.   This design facilitates transparent and more complete validation of input data with resulting benefits of application robustness to end users. 

\paragraph{Compile-time to runtime variants} as seen in Tables~\ref{tab:GitHubActions} and ~\ref{tab:self-hosted}, QMCPACK has a large number of compile-time configuration options that narrows the applicability of the resulting binary to specific cases.
These were introduced early in the development of the application.
With the availability of rich meta-programming capabilities via C\texttt{++} templates, we aim to migrate
some of these compile-time configurations to runtime variants.
Thus, a single executable can support several of these configurations without the need to recompile or distribute several versions of the code.
The most critical case is the unification of the use of complex and real numbers in wavefunction calculations.
This is an ongoing effort that would simplify use of the application, reduce testing burden, and simplify installation/distribution.
